% Generated by scripts/gen_blueprint.py from blueprint/nodes/09-baker.toml. Do not edit.

\chapter{Baker's theorem, by Schneider–Lang in several variables}
\label{chap:09-baker}

Baker's theorem in its qualitative, inhomogeneous form (Theorem~\ref{thm:baker_linear_forms_in_logarithms}): if $\ell_1,\dots,\ell_n$ are logarithms of algebraic numbers, linearly independent over $\Q$, then $1,\ell_1,\dots,\ell_n$ are linearly independent over $\Qbar$. The route is the one Bertrand and Masser found \cite{BM80,Mas81}, as Waldschmidt writes it in Chapter~4 of \cite{Wal00}: Baker's theorem for a basis of a number field (the book's Theorem~4.5), from the criterion of Schneider--Lang for $\C^{d_0}\times(\C^{\times})^{d_1}$ with $d_0\le1$ (its Corollary~4.2). The criterion is proved directly, as in the book's \S4.6: an auxiliary function from Siegel's lemma, Liouville's inequality at the points of a grid, a Schwarz lemma for Cartesian products (the book's Proposition~4.7) at the first derivative that does not vanish, and a contradiction. No zero estimate is needed. The Schwarz lemma comes from one-variable Hermite division, one coordinate at a time.

Step~5 of the book's \S4.6 needed a repair: vanishing to a given order in each coordinate does not survive the change of variables that the proof makes, while vanishing to a given total order does; only the constants change. No new mathematics is claimed.

\textbf{Notation.} $\iota$ is a finite set, $n=|\iota|$ unless said otherwise, and $\C^{\iota}$ carries the sup norm, so that its closed balls are polydiscs. For $w,z\in\C^{\iota}$, $\langle w,z\rangle=\sum_\nu w_\nu z_\nu$. For a function $F$ on $\C^{\iota}$ and a list $L=(L(0),\dots,L(k-1))$ of $k$ coordinate directions,
\[ D_LF(z)=D^{k}F(z)\bigl(e_{L(0)},\dots,e_{L(k-1)}\bigr), \]
where $D^{k}F$ is the $k$-th Fréchet derivative and $e_\nu$ are the coordinate vectors; every derivative of $F$ of total order $k$ vanishes at $z$ when $D^{k}F(z)=0$. Given $x_1,\dots,x_{d_1}\in\C^{\iota}$, $k_0\in\iota$ and integers $T_0,T_1\ge0$, the exponential monomials are $z_{k_0}^{\tau}e^{\langle t_1x_1+\dots+t_{d_1}x_{d_1},\,z\rangle}$ for $0\le\tau\le T_0$ and $t\in\{0,\dots,T_1\}^{d_1}$, and for coefficients $p=(p_{\tau,t})$
\[ F_p(z)=\sum_{\tau,t}p_{\tau,t}\,z_{k_0}^{\tau}\,e^{\langle t_1x_1+\dots+t_{d_1}x_{d_1},\,z\rangle}. \]
The grid points of a family $(y_j)_{j\in\iota}$ in $\C^{\iota}$ are the points $\sum_js_jy_j$ with integers $0\le s_j<S_1$.

\section{The Cartesian Schwarz lemma}

% Prove2Me: Transcendence.partials_eq_iteratedFDeriv -- https://prove2.me/theorems/a03f7908-0087-4475-9655-737ef73c454e
\begin{lemma}[Partial derivatives along a list of directions]
\label{lem:partials_eq_iteratedFDeriv}
\lean{Transcendence.partials_eq_iteratedFDeriv}
\leanok
Let $\mathbb{K}$ be a nontrivially normed field, $F$ a normed space over $\mathbb{K}$, and $g\colon\mathbb{K}^{\iota}\to F$ a function of class $C^{k}$. For $i\in\iota$ and a function $h$ on $\mathbb{K}^{\iota}$, let $\partial_ih(z)$ be the derivative at $w=z_i$ of $w\mapsto h(z[i:=w])$, where $z[i:=w]$ is $z$ with its coordinate $z_i$ replaced by $w$. Then for every list $L$ of $k$ coordinate directions and every $z\in\mathbb{K}^{\iota}$
\[ \partial_{L(0)}\,\partial_{L(1)}\cdots\partial_{L(k-1)}\,g(z)=D_Lg(z), \]
the partial derivative in the direction $L(k-1)$ being taken first.

{\small\emph{Source:} Standard.}
\end{lemma}

\begin{proof}
\leanok
\end{proof}

% Prove2Me: Transcendence.polydisc_cauchy -- https://prove2.me/theorems/623dbe11-9b39-48a3-aa5c-a35bb6a4613c
\begin{lemma}[Cauchy's inequalities on a polydisc]
\label{lem:polydisc_cauchy}
\lean{Transcendence.polydisc_cauchy}
\leanok
Let $F$ be an entire function on $\C^{\iota}$ with $|F|\le M$ on the closed polydisc of radius $\rho>0$ about $p$. For a list $L$ of $k$ coordinate directions, let $\sigma_\nu$ be the number of times $\nu$ occurs in $L$. Then
\[ |D_LF(p)|\le\Bigl(\prod_\nu\sigma_\nu!\Bigr)\frac{M}{\rho^{k}}. \]

{\small\emph{Source:} Classical; in the form of \cite[p.~XVII]{Wal00}.}
\end{lemma}

\begin{proof}
\leanok
\uses{lem:partials_eq_iteratedFDeriv}
Through Lemma~\ref{lem:partials_eq_iteratedFDeriv}, by Cauchy's inequality in one variable at a time.
\end{proof}

% Prove2Me: Transcendence.hermite_basis -- https://prove2.me/theorems/2e4bdbd5-ff19-4bd2-a497-f14f53276e86
\begin{lemma}[Hermite interpolation]
\label{lem:hermite_basis}
\lean{Transcendence.hermite_basis}
\leanok
Let $K$ be a field of characteristic $0$, $E\subset K$ a finite set and $S_0\in\N$, and write $q^{(k)}$ for the $k$-th formal derivative of $q\in K[X]$. There are polynomials $b_{\zeta,k}\in K[X]$ ($\zeta\in E$, $0\le k<S_0$) of degree less than $|E|S_0$ such that, for $\zeta'\in E$ and $0\le k'<S_0$, $b_{\zeta,k}^{(k')}(\zeta')$ is $1$ if $(\zeta',k')=(\zeta,k)$ and $0$ otherwise, and every $q\in K[X]$ of degree less than $|E|S_0$ satisfies
\[ q=\sum_{\zeta\in E}\ \sum_{k<S_0}q^{(k)}(\zeta)\,b_{\zeta,k}. \]

{\small\emph{Source:} Standard (Hermite interpolation).}
\end{lemma}

\begin{proof}
\leanok
\end{proof}

% Prove2Me: Transcendence.hermite_division_bound -- https://prove2.me/theorems/55157768-e9dd-4629-8ed1-e71c85e6e6da
\begin{lemma}[Hermite division with bounds]
\label{lem:hermite_division_bound}
\lean{Transcendence.hermite_division_bound}
\leanok
Let $0<r$ and $5r\le R$, let $Z$ be a finite multiset of complex numbers of modulus at most $r$, put $p=|Z|$ and $P(x)=\prod_{\zeta\in Z}(x-\zeta)$, and let $u$ be an entire function with $|u|\le M$ on the closed disc of radius $R$ about $0$. Then $u=\rho+Pq$ for a polynomial $\rho$ of degree less than $p$ and an entire function $q$, such that $\rho^{(k)}(\zeta)=u^{(k)}(\zeta)$ whenever $k$ is less than the multiplicity of $\zeta$ in $Z$, and for $|x|\le R$
\[ 2|\rho(x)|+M\le3^{p}M,\qquad|q(x)|\le(3/R)^{p}M. \]

{\small\emph{Source:} \cite[Lemma 4.8 c)--d), p.~123]{Wal00} in one variable, with the bounds of Step 2.4 of its proof (pp.~125--126).}
\end{lemma}

\begin{proof}
\leanok
Induction on the nodes, dividing by one factor $x-\zeta$ at a time, as in Step 2.4 of the book; $\rho$ and $q$ are the book's $f_0$ and $f_n$.
\end{proof}

% Prove2Me: Transcendence.coord_hermite_step -- https://prove2.me/theorems/88e63d37-3daa-47b5-a7c1-54ccbc4f22bc
\begin{lemma}[Hermite interpolation in one coordinate]
\label{lem:coord_hermite_step}
\lean{Transcendence.coord_hermite_step}
\leanok
Let $g$ be an entire function on $\C^{\iota}$ with $|g|\le K$ on the closed polydisc of radius $R$ about $0$, let $i\in\iota$ and $S_0\in\N$, and let $E$ be a finite set of complex numbers of modulus at most $r$, where $0<r$ and $5r\le R$; put $p=|E|S_0$. There is an entire function $h$ on $\C^{\iota}$ such that:
\begin{itemize}
\item $|h|\le3^{p}K$ on the closed polydisc of radius $R$ about $0$;
\item $|g-h|\le(2r)^{p}(3/R)^{p}K$ on the closed polydisc of radius $r$ about $0$;
\item for every point $\xi$ and every list $v$ of $m$ coordinate directions different from $i$: if $D_{(v,i^{k})}\,g(\xi[i:=\zeta])=0$ for every $\zeta\in E$ and every $k<S_0$, where $(v,i^{k})$ is the list $v$ followed by $k$ times $i$, then $D_vh(\xi)=0$.
\end{itemize}

{\small\emph{Source:} A step of the proof of \cite[Proposition 4.7, pp.~122--130]{Wal00}: the case $m=n$ of its Lemma 4.8 (pp.~123--126) in one coordinate, with the bounds of Step 2.4.}
\end{lemma}

\begin{proof}
\leanok
\uses{lem:hermite_division_bound,lem:hermite_basis,lem:partials_eq_iteratedFDeriv}
$h$ is the Hermite interpolant of $g$ in the variable $z_i$ at the nodes of $E$, each with multiplicity $S_0$: Lemma~\ref{lem:hermite_division_bound} on every slice, written in the fixed basis of Lemma~\ref{lem:hermite_basis}, so that the remainder depends on the other variables only through the jets of $g$ at the nodes. Only one-variable complex analysis is used.
\end{proof}

% Prove2Me: Transcendence.cartesian_schwarz -- https://prove2.me/theorems/ad262221-bdaa-4542-9af2-2228372ffe2c
\begin{theorem}[Schwarz's lemma for Cartesian products]
\label{thm:cartesian_schwarz}
\lean{Transcendence.cartesian_schwarz}
\leanok
Let $n\ge1$ and let $f$ be an entire function on $\C^{n}$. Let $E_1,\dots,E_n$ be sets of $S_1$ complex numbers each, all of modulus at most $r>0$, and let $R\ge5r$. Suppose that $|f|\le M$ on the closed polydisc of radius $R$ about $0$, and that every derivative of $f$ of total order less than $nS_0$ vanishes at every point of $E_1\times\dots\times E_n$. Then on the closed polydisc of radius $r$ about $0$
\[ |f|\le n\Bigl(\frac{2\cdot3^{n}r}{R}\Bigr)^{S_0S_1}M. \]

{\small\emph{Source:} \cite[Proposition 4.7, pp.~122--130]{Wal00}, in the form used in its \S4.6, with $2\cdot3^{n}$ in place of $18^{n}$.}
\end{theorem}

\begin{proof}
\leanok
\uses{lem:coord_hermite_step}
Lemma~\ref{lem:coord_hermite_step} in each coordinate in turn, telescoping the differences. The book assumes vanishing with every exponent below $S_0$; vanishing by total order, assumed here, is what the change of variables of \S4.6 preserves.
\end{proof}

\section{Siegel's lemma and the Taylor bounds}

% Prove2Me: Transcendence.thue_siegel_real -- https://prove2.me/theorems/4fe54bab-8cb3-4b35-aaa3-3e03bc6f3a39
\begin{lemma}[Thue–Siegel's lemma for real linear forms]
\label{lem:thue_siegel_real}
\lean{Transcendence.thue_siegel_real}
\leanok
Let $J$ and $\Lambda$ be finite sets, $\mu=|J|$ and $\nu=|\Lambda|$, and let $v_{j\lambda}\in\R$ with $\sum_\lambda|v_{j\lambda}|\le C$ for every $j$, where $C>0$. Let $X\ge0$ and $\ell\ge1$ be integers with $\ell^{\mu}<(X+1)^{\nu}$. Then there is $\xi\in\Z^{\Lambda}$, $\xi\neq0$, with $|\xi_\lambda|\le X$ for every $\lambda$ and
\[ \Bigl|\sum_\lambda v_{j\lambda}\xi_\lambda\Bigr|\le\frac{CX}{\ell}\qquad(j\in J). \]

{\small\emph{Source:} \cite[Lemma 4.11, pp.~132--133]{Wal00}, with any real $C>0$ in place of an integer.}
\end{lemma}

\begin{proof}
\leanok
Dirichlet's box principle.
\end{proof}

% Prove2Me: Transcendence.taylor_tail_le -- https://prove2.me/theorems/7815b36a-411f-427d-824a-1598c03f4e04
\begin{lemma}[The tail of a Taylor series]
\label{lem:taylor_tail_le}
\lean{Transcendence.taylor_tail_le}
\leanok
Let $E$ be a complex normed space and $G\colon E\to\C$ analytic at every point, with $|G|\le M$ on the closed ball of radius $R$ about $0$, and let $0<r<R$. For every $z\in E$ with $\|z\|\le r$ and every $T\in\N$
\[ \Bigl|G(z)-\sum_{k<T}\frac{1}{k!}\,D^{k}G(0)(z,\dots,z)\Bigr|\le\frac{M(r/R)^{T}}{1-r/R}. \]

{\small\emph{Source:} The tail bound in the proof of \cite[Lemma 4.13, pp.~134--135]{Wal00}, with $1/(1-r/R)$ in place of $1+\sqrt{T}$.}
\end{lemma}

\begin{proof}
\leanok
Cauchy's inequality on the line through $z$ bounds each term, and the terms are summed as a geometric series.
\end{proof}

% Prove2Me: Transcendence.taylor_coeff_forms -- https://prove2.me/theorems/545622c4-c143-48c2-a85d-280d5d1be38f
\begin{lemma}[Taylor coefficients as linear forms]
\label{lem:taylor_coeff_forms}
\lean{Transcendence.taylor_coeff_forms}
\leanok
Let $\Lambda$ be a finite set, let $\varphi_\lambda$ ($\lambda\in\Lambda$) be entire functions on $\C^{\iota}$ with $|\varphi_\lambda|\le B_\lambda$ on the closed polydisc of radius $r>0$ about $0$, and let $T\in\N$. There are $u_{\tau\lambda}\in\C$, for $\lambda\in\Lambda$ and $\tau\in\{0,\dots,T-1\}^{\iota}$, with $|u_{\tau\lambda}|\le B_\lambda$, such that for all $c\in\C^{\Lambda}$ and $z\in\C^{\iota}$ the function $F=\sum_\lambda c_\lambda\varphi_\lambda$ satisfies
\[ \sum_{k<T}\frac{1}{k!}\,D^{k}F(0)(z,\dots,z)=\sum_{\tau}\Bigl(\sum_\lambda u_{\tau\lambda}c_\lambda\Bigr)\prod_{\nu\in\iota}\Bigl(\frac{z_\nu}{r}\Bigr)^{\tau_\nu}. \]

{\small\emph{Source:} A step of the proof of \cite[Proposition 4.10, p.~135]{Wal00}, with Cauchy's inequalities (p.~XVII).}
\end{lemma}

\begin{proof}
\leanok
\uses{lem:polydisc_cauchy}
Expand the Taylor polynomial in monomials; Lemma~\ref{lem:polydisc_cauchy} bounds the coefficients.
\end{proof}

% Prove2Me: Transcendence.siegel_small_values_of_count -- https://prove2.me/theorems/1949ad28-451d-4549-899e-5b22a33f506c
\begin{lemma}[A small integer combination, parameters free]
\label{lem:siegel_small_values_of_count}
\lean{Transcendence.siegel_small_values_of_count}
\leanok
Let $\Lambda$ be a finite set, $L=|\Lambda|$, and let $\varphi_\lambda$ ($\lambda\in\Lambda$) be entire functions on $\C^{\iota}$ with $|\varphi_\lambda|\le B_\lambda$ on the closed polydisc of radius $R$ about $0$ and $\sum_\lambda B_\lambda\le C$, where $C>0$. Let $0<r<R$, and let $T$, $X$ and $\ell\ge1$ be integers with $\ell^{2T^{n}}<(X+1)^{L}$. Then there are integers $p_\lambda$, not all zero, with $|p_\lambda|\le X$, such that on the closed polydisc of radius $r$ about $0$
\[ \Bigl|\sum_\lambda p_\lambda\varphi_\lambda\Bigr|\le T^{n}\cdot\frac{2CX}{\ell}+\frac{XC(r/R)^{T}}{1-r/R}. \]

{\small\emph{Source:} \cite[\S4.5: Lemmas 4.11--4.13 and the proof of Proposition 4.10, pp.~132--136]{Wal00}, with $T$, $X$ and $\ell$ left free.}
\end{lemma}

\begin{proof}
\leanok
\uses{lem:thue_siegel_real,lem:taylor_tail_le,lem:taylor_coeff_forms}
Lemma~\ref{lem:thue_siegel_real}, applied to the real and imaginary parts of the forms of Lemma~\ref{lem:taylor_coeff_forms}, makes the Taylor polynomial small; Lemma~\ref{lem:taylor_tail_le} bounds the tail.
\end{proof}

% Prove2Me: Transcendence.siegel_small_values_parameters -- https://prove2.me/theorems/65e8b744-5d79-469e-aee5-891c31c9b115
\begin{lemma}[The parameters of Proposition 4.10]
\label{lem:siegel_small_values_parameters}
\lean{Transcendence.siegel_small_values_parameters}
\leanok
Let $n\ge1$ and $L$ be integers and $N,U,V,r,R$ real numbers with $r>0$, $W=N+U+V\ge12n^{2}$, $er\le R\le re^{W/6}$ and
\[ (2W)^{n+1}\le L\,N\,\bigl(\log(R/r)\bigr)^{n}. \]
Then there are integers $T$, $X$ and $\ell\ge1$ with $X\le e^{N}$ and $\ell^{2T^{n}}<(X+1)^{L}$ such that
\[ T^{n}\cdot\frac{2e^{U}X}{\ell}\le\frac34e^{-V},\qquad\frac{Xe^{U}(r/R)^{T}}{1-r/R}\le\frac14e^{-V}. \]

{\small\emph{Source:} The numerical part of the proof of \cite[Proposition 4.10, pp.~135--136]{Wal00}.}
\end{lemma}

\begin{proof}
\leanok
The book's choices: $\frac43W\le T\log(R/r)<\frac43W+\log(R/r)$, $X=\lfloor e^{N}\rfloor$, and $\ell=\lceil\frac83T^{n}e^{W}\rceil$ cells for the box principle.
\end{proof}

% Prove2Me: Transcendence.siegel_small_values -- https://prove2.me/theorems/22592b42-9f29-44cd-a4a9-45463409e037
\begin{proposition}[Thue–Siegel: a small auxiliary function]
\label{prop:siegel_small_values}
\lean{Transcendence.siegel_small_values}
\leanok
Let $n=|\iota|\ge1$, let $\Lambda$ be a finite set, and let $\varphi_\lambda$ ($\lambda\in\Lambda$) be entire functions on $\C^{\iota}$. Let $N,U,V,r>0$ and $R$ be real numbers with $W=N+U+V\ge12n^{2}$ and $er\le R\le re^{W/6}$. Suppose that $|\varphi_\lambda|\le B_\lambda$ on the closed polydisc of radius $R$ about $0$, with $\sum_\lambda B_\lambda\le e^{U}$, and that
\[ (2W)^{n+1}\le|\Lambda|\,N\,\bigl(\log(R/r)\bigr)^{n}. \]
Then there are integers $p_\lambda$, not all zero, with $|p_\lambda|\le e^{N}$, such that $\bigl|\sum_\lambda p_\lambda\varphi_\lambda\bigr|\le e^{-V}$ on the closed polydisc of radius $r$ about $0$.

{\small\emph{Source:} \cite[Proposition 4.10 with Lemmas 4.11--4.13, pp.~131--136]{Wal00}.}
\end{proposition}

\begin{proof}
\leanok
\uses{lem:siegel_small_values_of_count,lem:siegel_small_values_parameters}
Lemma~\ref{lem:siegel_small_values_of_count} with the parameters of Lemma~\ref{lem:siegel_small_values_parameters} and $C=e^{U}$.
\end{proof}

\section{Exponential monomials}

% Prove2Me: Transcendence.exp_monomial_derivs -- https://prove2.me/theorems/d2ef3d38-bf65-41b0-86fd-9603c916e1a3
\begin{lemma}[Derivatives of one exponential monomial]
\label{lem:exp_monomial_derivs}
\lean{Transcendence.exp_monomial_derivs}
\leanok
Let $K$ be a number field with an embedding $\varphi\colon K\to\C$, let $w\in K^{\iota}$, $\tau\in\N$, $k_0\in\iota$, and $g(z)=z_{k_0}^{\tau}\exp\bigl(\sum_\nu\varphi(w_\nu)z_\nu\bigr)$. Let $q\in\C^{\iota}$ and $v\in K$, with $\varphi(v)=q_{k_0}$ if $\tau>0$. Suppose that $A\ge1$, $\houseof{w_\nu}\le A$ for every $\nu$, $\houseof{v}\le B$, and that $\delta w_\nu$ and $\delta v$ are algebraic integers for an integer $\delta$. Then for every list $L$ of $k$ coordinate directions
\[ D_Lg(q)=\varphi(\gamma)\exp\Bigl(\sum_\nu\varphi(w_\nu)q_\nu\Bigr) \]
for some $\gamma\in K$ with $\houseof{\gamma}\le A^{k}(B+k)^{\tau}$ and $\delta^{k+\tau}\gamma$ an algebraic integer.

{\small\emph{Source:} \cite[Lemma 4.9, p.~130]{Wal00}, in the form used in its \S4.6.}
\end{lemma}

\begin{proof}
\leanok
\end{proof}

% Prove2Me: Transcendence.exp_monomials_ne_zero -- https://prove2.me/theorems/1841b1c4-1bc4-4659-9c82-bb3bf5723fb0
\begin{lemma}[Exponential monomials are independent]
\label{lem:exp_monomials_ne_zero}
\lean{Transcendence.exp_monomials_ne_zero}
\leanok
Let $x_1,\dots,x_{d_1}\in\C^{\iota}$ be linearly independent over $\Q$, let $k_0\in\iota$ and $T_0,T_1\in\N$, and let the coefficients $p_{\tau,t}\in\C$ be not all zero. Then $D^{k}F_p(0)\neq0$ for some $k$.

{\small\emph{Source:} \cite[Exercises 2.4--2.5, p.~60]{Wal00}, in the form used in its \S4.6.}
\end{lemma}

\begin{proof}
\leanok
\uses{lem:expPoly_ne_zero}
The frequencies $\sum_it_ix_i$ are pairwise distinct. On a suitable line, $F_p$ is a one-variable exponential polynomial with distinct frequencies, which is not identically zero by Lemma~\ref{lem:expPoly_ne_zero}.
\end{proof}

% Prove2Me: Transcendence.exp_monomials_deriv_lattice -- https://prove2.me/theorems/6e37a083-ceb0-42a9-aca6-1c6af7b701b4
\begin{lemma}[Derivatives at the grid points]
\label{lem:exp_monomials_deriv_lattice}
\lean{Transcendence.exp_monomials_deriv_lattice}
\leanok
Let $K$ be a number field with an embedding $\varphi\colon K\to\C$, let $k_0\in\iota$, and let $x_1,\dots,x_{d_1}$ and $y_j$ ($j\in\iota$) be vectors in $\C^{\iota}$. Suppose that $x_{i\nu}=\varphi(\xi_{i\nu})$ and $e^{\langle x_i,y_j\rangle}=\varphi(a_{ij})$ with $\xi_{i\nu},a_{ij}\in K$, and that $\eta_j\in K$ satisfy $\varphi(\eta_j)=y_{jk_0}$ if $T_0>0$. Suppose that an integer $\delta$ makes every $\delta\xi_{i\nu}$, $\delta\eta_j$ and $\delta a_{ij}$ an algebraic integer, and that all their houses are at most $H\ge1$. Let $T_1\ge1$ and $X\in\R$, and let the coefficients $p_{\tau,t}$ be integers with $|p_{\tau,t}|\le X$. Then at every grid point $\sum_js_jy_j$ and for every list $L$ of $k$ coordinate directions, $D_LF_p\bigl(\sum_js_jy_j\bigr)=\varphi(\gamma)$ for some $\gamma\in K$ with
\[ \houseof{\gamma}\le(T_0+1)(T_1+1)^{d_1}X\,\bigl((d_1+1)T_1H\bigr)^{k}\,(nS_1H+k)^{T_0}\,H^{d_1nT_1S_1} \]
and $\delta^{k+T_0+d_1nT_1S_1}\gamma$ an algebraic integer.

{\small\emph{Source:} A step of \cite[\S4.6, step 2, p.~137]{Wal00}, with Lemma 4.9 (pp.~130--131).}
\end{lemma}

\begin{proof}
\leanok
\uses{lem:exp_monomial_derivs}
Lemma~\ref{lem:exp_monomial_derivs} for each monomial, at the grid point, where $e^{\langle\sum_it_ix_i,\,\sum_js_jy_j\rangle}=\varphi\bigl(\prod_{i,j}a_{ij}^{t_is_j}\bigr)$.
\end{proof}

% Prove2Me: Transcendence.exp_monomials_liouville_lower -- https://prove2.me/theorems/3e0a1030-aedc-4ace-a12b-8c4476667212
\begin{lemma}[Liouville's lower bound at the grid points]
\label{lem:exp_monomials_liouville_lower}
\lean{Transcendence.exp_monomials_liouville_lower}
\leanok
Let $k_0\in\iota$ and $d_0\le1$, and let $x_1,\dots,x_{d_1}$ and $y_j$ ($j\in\iota$) be vectors in $\C^{\iota}$ such that every coordinate $x_{i\nu}$ and every $e^{\langle x_i,y_j\rangle}$ is algebraic and, if $d_0=1$, every $y_{jk_0}$ is algebraic. There is $C\ge1$ with the following property. Let $T\ge1$ and $S_1$ be integers and $N\ge0$, and let $F_p$, with $T_0=d_0T$ and $T_1=T$, have integer coefficients with $|p_{\tau,t}|\le e^{N}$. If $\Delta=D_LF_p\bigl(\sum_js_jy_j\bigr)\neq0$ for a list $L$ of $k$ coordinate directions and a grid point $\sum_js_jy_j$, then
\[ \log|\Delta|\ge-C\bigl(N+k(1+\log T)+T(S_1+\log(1+k))\bigr). \]

{\small\emph{Source:} \cite[(4.14), \S4.6, step 2, p.~137]{Wal00}, with a bracket of a slightly different shape.}
\end{lemma}

\begin{proof}
\leanok
\uses{lem:exp_monomials_deriv_lattice,lem:liouville_house}
A number field containing the data, a common denominator and a bound for the houses are built inside. Lemma~\ref{lem:exp_monomials_deriv_lattice} writes $\Delta$ as an algebraic number, and Liouville's inequality (Lemma~\ref{lem:liouville_house}) bounds it from below.
\end{proof}

% Prove2Me: Transcendence.exists_exp_monomials_small_on_grid -- https://prove2.me/theorems/f435be21-31ad-48d9-a07d-c81bc63cb88c
\begin{lemma}[An auxiliary function small at the grid points]
\label{lem:exists_exp_monomials_small_on_grid}
\lean{Transcendence.exists_exp_monomials_small_on_grid}
\leanok
Let $k_0\in\iota$, $d_0\le1$ and $d=d_0+d_1$, and let $x_1,\dots,x_{d_1}$ and $y_j$ ($j\in\iota$) be vectors in $\C^{\iota}$ with $\sum_{i,\nu}|x_{i\nu}|\le c_x$ and $\sum_j\|y_j\|\le c_y$. Let $S_1,T\ge1$ be integers and $U,N>0$ and $E$ real numbers; put $r=(c_y+2)S_1$ and $R=Er$, and assume
\[ 12n^{2}\le N+2U,\qquad e\le E\le e^{(N+2U)/6}, \]
\[ (T+1)^{d}R^{T}e^{c_xTR}\le e^{U},\qquad\bigl(2(N+2U)\bigr)^{n+1}\le(T+1)^{d}N(\log E)^{n}. \]
Then there are integers $p_{\tau,t}$, not all zero, with $|p_{\tau,t}|\le e^{N}$, such that, with $T_0=d_0T$ and $T_1=T$:
\begin{itemize}
\item $\bigl|D_LF_p\bigl(\sum_js_jy_j\bigr)\bigr|\le k!\,e^{-U}$ at every grid point and for every list $L$ of $k$ coordinate directions;
\item $|F_p(w)|\le(T+1)^{d}e^{N}\rho^{T}e^{c_xT\rho}$ whenever $\rho\ge1$ and $\|w\|\le\rho$.
\end{itemize}

{\small\emph{Source:} \cite[\S4.6, step 3 and (4.16), pp.~138--139]{Wal00}.}
\end{lemma}

\begin{proof}
\leanok
\uses{prop:siegel_small_values,lem:polydisc_cauchy}
Proposition~\ref{prop:siegel_small_values} for the $(T+1)^{d}$ monomials, with $V=U$ and the radii $r$ and $R$; Cauchy's inequalities (Lemma~\ref{lem:polydisc_cauchy}) on the unit polydisc about each grid point give the first bound.
\end{proof}

\section{The criterion of Schneider–Lang}

% Prove2Me: Transcendence.grid_schwarz_cauchy -- https://prove2.me/theorems/8b855f06-5713-4fa8-9781-42d2ffa24c05
\begin{lemma}[Schwarz's lemma at the grid points, then Cauchy]
\label{lem:grid_schwarz_cauchy}
\lean{Transcendence.grid_schwarz_cauchy}
\leanok
Let $n=|\iota|\ge1$ and let $(y_j)_{j\in\iota}$ be a basis of $\C^{\iota}$. There is $c\ge1$ with the following property. Let $F$ be an entire function on $\C^{\iota}$, let $S_0$, $S_1\ge1$ and $M\ge nS_0$ be integers, and let $\rho\ge1$ and $B$ be real numbers. Suppose that every derivative of $F$ of total order less than $M$ vanishes at every grid point, and that $|F(w)|\le B$ whenever $\|w\|\le cS_1\rho$. Then at every grid point and for every list $L$ of $M$ coordinate directions
\[ \Bigl|D_LF\Bigl(\sum_js_jy_j\Bigr)\Bigr|\le M!\cdot n\,\rho^{-S_0S_1}B. \]

{\small\emph{Source:} \cite[\S4.6, step 5, p.~140]{Wal00}, with Proposition 4.7 (p.~122) and Cauchy's inequalities (p.~XVII).}
\end{lemma}

\begin{proof}
\leanok
\uses{thm:cartesian_schwarz,lem:polydisc_cauchy}
Theorem~\ref{thm:cartesian_schwarz} for $F\bigl(\sum_jz_jy_j\bigr)$ on the grid $\{0,\dots,S_1-1\}^{n}$, then Lemma~\ref{lem:polydisc_cauchy} on the unit polydisc about the grid point. This is where the book's vanishing hypothesis, with every exponent below $S_0$, is not preserved by the change of variables $z\mapsto\sum_jz_jy_j$; vanishing by total order is.
\end{proof}

% Prove2Me: Transcendence.schneider_lang_vanishing_ineq -- https://prove2.me/theorems/36ba44ee-12ee-466a-9ed0-d6d75c56e741
\begin{lemma}[The inequality behind (4.17)]
\label{lem:schneider_lang_vanishing_ineq}
\lean{Transcendence.schneider_lang_vanishing_ineq}
\leanok
Let $n\ge1$ be an integer and $C\ge1$. There is an integer $s_0$ such that for every integer $S_1\ge s_0$ there is $e_0$ with the following property. Let $T\ge1$ and $E\ge e_0$ be integers with $\log T\le n\log S_1+n\log E$, and put
\[ U=\frac{S_1ET\log E}{2C\cdot6^{n+1}},\qquad N=\frac{U}{2C}. \]
Then for every integer $k<nET$
\[ C\bigl(N+k(1+\log T)+T(S_1+\log(1+k))\bigr)+\log k!<U. \]

{\small\emph{Source:} \cite[\S4.6, step 4 and (4.17), p.~139]{Wal00}.}
\end{lemma}

\begin{proof}
\leanok
\end{proof}

% Prove2Me: Transcendence.schneider_lang_extrapolation_ineq -- https://prove2.me/theorems/5b268b48-6fc2-4f44-80f2-80ac1a58586e
\begin{lemma}[The inequality behind (4.19)]
\label{lem:schneider_lang_extrapolation_ineq}
\lean{Transcendence.schneider_lang_extrapolation_ineq}
\leanok
Let $n\ge1$ and $d$ be integers, $C\ge1$, $c_x\in\R$ and $c_F\ge1$. There is an integer $s_0$ such that for every integer $S_1\ge s_0$ there is $e_0$ with the following property. Let $T\ge1$ and $E\ge e_0$ be integers with $\log T\le n\log S_1+n\log E$, and let $U$ and $N$ be as in Lemma~\ref{lem:schneider_lang_vanishing_ineq}. Then for all integers $u\ge ET$ and $M\le2nu$, with $\rho'=c_FS_1u/T$,
\[ \begin{aligned} &C\bigl(N+M(1+\log T)+T(S_1+\log(1+M))\bigr)+\log M!+\log n\\ &\qquad+d\log(T+1)+N+T\log\rho'+c_xT\rho'<uS_1\log(u/T). \end{aligned} \]

{\small\emph{Source:} \cite[\S4.6, step 6 and (4.19), pp.~140--141]{Wal00}.}
\end{lemma}

\begin{proof}
\leanok
The book gives these asymptotics without proof; here they hold uniformly in $u\ge ET$, the book's $S_0'$.
\end{proof}

% Prove2Me: Transcendence.exists_schneider_lang_parameters -- https://prove2.me/theorems/483ae7ac-c47e-455f-9cdd-bcd2d8128133
\begin{lemma}[The parameters of the criterion]
\label{lem:exists_schneider_lang_parameters}
\lean{Transcendence.exists_schneider_lang_parameters}
\leanok
Let $1\le n<d$ be integers, $C\ge1$, $c_x,c_y\ge0$ and $c_F\ge1$. There are integers $S_1,T,E\ge1$ and real numbers $U,N>0$ such that, with $r=(c_y+2)S_1$ and $R=Er$:
\begin{enumerate}
\item the four assumptions of Lemma~\ref{lem:exists_exp_monomials_small_on_grid} hold;
\item the inequality of Lemma~\ref{lem:schneider_lang_vanishing_ineq} holds for every integer $k<nET$;
\item the inequality of Lemma~\ref{lem:schneider_lang_extrapolation_ineq} holds for all integers $u\ge ET$ and $M\le2nu$, with $\rho'=c_FS_1u/T$.
\end{enumerate}

{\small\emph{Source:} \cite[\S4.6: the conditions (4.15)--(4.19) and the choice of parameters in step 6, pp.~138--141]{Wal00}.}
\end{lemma}

\begin{proof}
\leanok
\uses{lem:schneider_lang_vanishing_ineq,lem:schneider_lang_extrapolation_ineq}
As in the book: $S_1$ fixed and large, then $T=(S_1m)^{n}$ and $E=S_1^{d-1-n}m^{d-n}$ for a large integer $m$, $U=S_1ET\log E/(2C\cdot6^{n+1})$ and $N=U/(2C)$. Lemmas~\ref{lem:schneider_lang_vanishing_ineq} and~\ref{lem:schneider_lang_extrapolation_ineq} give conditions 2 and 3.
\end{proof}

% Prove2Me: Transcendence.schneider_lang_cartesian -- https://prove2.me/theorems/0d6daaa2-4460-46f7-a092-0052518d8136
\begin{theorem}[The criterion of Schneider–Lang]
\label{thm:schneider_lang_cartesian}
\lean{Transcendence.schneider_lang_cartesian}
\leanok
Let $x_1,\dots,x_{d_1}\in\C^{\iota}$ have algebraic coordinates and be linearly independent over $\Q$, and let $(y_j)_{j\in\iota}$ be a basis of $\C^{\iota}$. Either let $d_0=0$, or let $d_0=1$ and fix a coordinate $k$ such that every $y_{jk}$ is algebraic. If $|\iota|<d_0+d_1$, then the numbers
\[ e^{\langle x_i,y_j\rangle}\qquad(1\le i\le d_1,\ j\in\iota) \]
are not all algebraic.

{\small\emph{Source:} Schneider (1949), Lang \cite{Lang66}; this form is \cite[Corollary 4.2, p.~117]{Wal00} for $d_0\le1$, proved directly as in its \S4.6 (pp.~136--141), on the route of Bertrand--Masser \cite{BM80} and Masser \cite{Mas81}.}
\end{theorem}

\begin{proof}
\leanok
\uses{lem:exp_monomials_liouville_lower,lem:exists_exp_monomials_small_on_grid,lem:grid_schwarz_cauchy,lem:exists_schneider_lang_parameters,lem:exp_monomials_ne_zero}
Suppose they are all algebraic. Lemma~\ref{lem:exists_schneider_lang_parameters} chooses the parameters, Lemma~\ref{lem:exists_exp_monomials_small_on_grid} an auxiliary function $F_p$, and Lemma~\ref{lem:exp_monomials_ne_zero} shows that it is not identically zero. By Lemma~\ref{lem:exp_monomials_liouville_lower} and condition 2, its derivatives of order less than $nET$ vanish at the grid points. At the first order at which one does not, Lemma~\ref{lem:grid_schwarz_cauchy} bounds it from above and Lemma~\ref{lem:exp_monomials_liouville_lower} from below, against condition 3.
\end{proof}

\section{Baker's theorem}

% Prove2Me: Transcendence.baker_number_field_basis -- https://prove2.me/theorems/46da3787-b843-4ab8-8c3d-1c627a7847aa
\begin{theorem}[Baker's theorem for a basis of a number field]
\label{thm:baker_number_field_basis}
\lean{Transcendence.baker_number_field_basis}
\leanok
Let $K\subset\C$ be a number field, $(\beta_k)$ a basis of $K$ over $\Q$, and $\ell_k$ logarithms of algebraic numbers. If $\sum_k\beta_k\ell_k$ is algebraic, then $\ell_k=0$ for every $k$.

{\small\emph{Source:} \cite[Theorem 4.5 with Lemma 4.6, pp.~119--121]{Wal00}; the route is that of Bertrand--Masser \cite{BM80} and Masser \cite{Mas81}.}
\end{theorem}

\begin{proof}
\leanok
\uses{lem:exp_ratMul_isAlgebraic,thm:schneider_lang_cartesian}
The matrix $(\sigma(\beta_k))$ of the embeddings $\sigma$ of $K$ is invertible, the trace form being non-degenerate. With $\lambda_\sigma=\sum_k\sigma(\beta_k)\ell_k$, Theorem~\ref{thm:schneider_lang_cartesian} is applied on the coordinates where $\lambda_\sigma\neq0$: with $d_0=1$ when $\lambda_\sigma\neq0$ at the inclusion $K\subset\C$, and with $d_0=0$ otherwise.
\end{proof}

% Prove2Me: Schanuel.baker_linear_forms_in_logarithms -- https://prove2.me/theorems/a0417130-5d91-4638-a293-bea5fb28fcf6
\begin{theorem}[Baker's theorem]
\label{thm:baker_linear_forms_in_logarithms}
\lean{Transcendence.baker_linear_forms_in_logarithms}
\leanok
Let $\ell_1,\dots,\ell_n$ be logarithms of algebraic numbers, linearly independent over $\Q$, and let $\beta_0,\beta_1,\dots,\beta_n$ be algebraic numbers, not all zero. Then
\[ \beta_0+\beta_1\ell_1+\dots+\beta_n\ell_n\neq0. \]
Equivalently, $1,\ell_1,\dots,\ell_n$ are linearly independent over $\Qbar$.

{\small\emph{Source:} Baker \cite{Bak66}; see \cite[Theorem 2.1]{Bak75}. The route is that of Bertrand--Masser \cite{BM80}, as in \cite[Chapter~4]{Wal00}.}
\end{theorem}

\begin{proof}
\leanok
\uses{lem:exp_ratMul_isAlgebraic,thm:baker_number_field_basis}
Write $\beta_i=\sum_kc_{ik}b_k$ in a basis $(b_k)$ of the number field generated by the $\beta_i$, with $c_{ik}\in\Q$. Then $\sum_kb_kL_k=-\beta_0$ is algebraic, where the $L_k=\sum_ic_{ik}\ell_i$ are logarithms of algebraic numbers (Lemma~\ref{lem:exp_ratMul_isAlgebraic}), so every $L_k$ is $0$ by Theorem~\ref{thm:baker_number_field_basis}. Independence over $\Q$ gives $c_{ik}=0$, hence $\beta_1=\dots=\beta_n=0$, and then $\beta_0=0$.
\end{proof}

% Prove2Me: DiazModulus.baker_two_logs -- https://prove2.me/theorems/f04bc443-81b3-49dd-a72f-1cbbcdf9b7b2
\begin{corollary}[Baker's theorem for two logarithms]
\label{cor:baker_two_logs}
\lean{Diaz.baker_two_logs}
\leanok
Let $x$ and $y$ be logarithms of algebraic numbers, linearly independent over $\Q$, and let $a,b$ be algebraic numbers, not both zero. Then $ax+by$ is transcendental.

{\small\emph{Source:} Baker \cite{Bak66}; the proof follows \cite[Chapter~4]{Wal00}, after Bertrand--Masser \cite{BM80} and Masser \cite{Mas81}.}
\end{corollary}

\begin{proof}
\leanok
\uses{thm:baker_linear_forms_in_logarithms}
If $ax+by=\beta$ were algebraic, Theorem~\ref{thm:baker_linear_forms_in_logarithms} with $\beta_0=-\beta$, $\beta_1=a$ and $\beta_2=b$ would fail.
\end{proof}
